\documentclass[a4paper]{article}

\usepackage{hyperref}
\hypersetup{
    colorlinks=true,
    linkcolor=blue,
    filecolor=magenta,      
    urlcolor=cyan,
}
\usepackage[utf8]{inputenc}
\usepackage[T1]{fontenc}
\usepackage{textcomp}
\usepackage[albanian]{babel}
\usepackage{amsmath, amssymb}
\usepackage{physics}
\usepackage{graphicx}
\usepackage{wrapfig}
\usepackage{lipsum}
\usepackage{biblatex}
\usepackage{float}

\newcommand{\RNum}[1]{\uppercase\expandafter{\romannumeral #1\relax}}

\title{Derivatet parciale dhe zbatimi i tyre në jetën e përditshme}
\author{Redon Tahiri\\Valdrin Ejupi\\Osman Bytyqi\\Valdete Salihi\\Rinesa Zogjani\\Samir Raci}
\date{Prill, 2020}

\begin{document}
\vbox{
	\centering
	\includegraphics[width=0.7\textwidth]{uplogo.png}
	\maketitle
}

\maketitle
\newpage
\thispagestyle{empty}
\begin{abstract}
	Ky punim është zgjedhur nga autorët si një projekt për lëndën ``Matematika \RNum{2}'' nga Universiteti i Prishtinës, Fakulteti i Inxhinierisë Elektrike dhe Kompjuterike. Qëllimi i këtij punimi është që të prezantojë derivatet parciale dhe përdorimin e tyre në matematikë dhe jetën e përditshme për lexuesit. Ky punim shkurtimisht i përkufizon derivatet parciale prej rendit të parë deri në ato të rendit të tretë, e përshkruan mënyrën se si ato mund të përdoren për optimizimin e funksioneve me shumë variabla, si dhe jep shembuj duke e përshkruar punën e bërë detajisht.
\end{abstract}
\newpage

\tableofcontents
\newpage

\section{Hyrje}
Në matematikë, një derivat parcial i një funksioni me dy ose më shumë variabla është derivati i tijë sipas njërës variabël kurse variablat e tjera mbahen konstante. Derivati parcial i një funksioni $f(x,y,\ldots)$ sipas variablës x shënohet me njërën nga këto mënyra:
\begin{center}$f'_x,f_x,\partial_xf, D_xf,\pdv{}{x}f$ ose $\pdv{f}{x}$\end{center}
Simboli i cili paraqet derivate parciale është $\partial$.

\subsection{Përkufizimi i derivateve parciale}
Sikurse derivatet e zakonshme, edhe derivatet parciale janë të definuara si një limit. Le të jetë $z=f(x,y)$. Shtesa parciale të funksionit $z=f(x,y)$ në lidhje me variablat e pavarura x,y i quajmë shprehjet:
\begin{align}
	\Delta z_x = f(x+\Delta x,y) - f(x,y) \text{ sipas x-it} \\
	\Delta z_y = f(x,y + \Delta y) - f(x,y) \text{ sipas y-it}.
\end{align}

Derivate parciale të funksionit $z = f(x,y)$ në pikën P(x,y) sipas variablave të pavaruara x respektivisht y i quajmë shprehjet:

\begin{gather}
	\pdv{z}{x} = \lim_{\Delta x\to 0}\frac{\Delta z_x}{\Delta x} = \lim_{\Delta x \to 0}\frac{f\left( x + \Delta x,y \right) - f(x,y) }{\Delta x}\\
	\pdv{z}{y} = \lim_{\Delta y\to 0}\frac{\Delta z_y}{\Delta y} = \lim_{\Delta y \to  0}\frac{f(x,y + \Delta y) - f(x,y)}{\Delta y}
\end{gather}

Në rastin e parë y konsiderohet konstantë, kurse në rastin e dytë, x konsiderohet konstantë.

\subsection{Derivatet Parciale të funksioneve të përbëra}
Nga funksionet me një variabël e dijmë që në qoftse $y=f(x)$ dhe $x=g(t)$, atëherë
\begin{gather*}
	\dv{y}{t} = \dv{y}{x} \dv{x}{t}
\end{gather*}
Në funksionet me dy ose më shumë variabla nëse kemi
\begin{gather*}
	z = f(x,y), x=g(t), y=h(t), \text{atëherë, }\dv{z}{t} \text{ do të jetë, }\\
	\dv{z}{t} = \pdv{f}{x}\dv{x}{t} + \pdv{f}{y}\dv{y}{t}
\end{gather*}
Këtu ne e diferencojmë f sipas të gjithave variablave në të dhë pastaj e shumzojmë secilën prej tyre nga derivati i asaj variable sipas t. Hapi perfundimtar është që ti mbledhim të gjithë këto së bashku.
\subsection{Shembull (1)}
Të gjindet derivati parcial i funksionit
\begin{equation}
	\label{label1}
	f(x,y)=x^4+6\sqrt{y}-10
\end{equation}
Së pari do ta gjejmë derivatin sipas x-it, kështu që variabla y do të trajtohet si konstantë. Derivati parcial i funksionit f sipas x do të jetë:
\begin{align*}
	f_x(x,y) = 4x^3`
\end{align*}
Vërejmë se termi i dytë dhe i tretë i funksionit diferencohen në zero në këtë rast. Kjo është sepse ne trajtuam variablën y si konstantë, dhe derivati i konstanës është zero.
Tani do ta gjejmë derivatin parcial të funksionit f sipas y-it, në këtë rast variabla x do të trajtohet si konstantë.
\begin{align*}
	f_y(x,y) = \frac{3}{\sqrt{y}}
\end{align*}

\section{Derivatet parciale të rendeve të larta}
Derivatet parciale e rendit të dytë ose më të lartë definohen në mënyrë analoge si derivatet e rendeve të larta të funksioneve me një variablël. Për funksionin $f(x,y,\ldots)$ derivati parcial i rendit të dytë i tij sipas x është thjeshtë derivati parcial i derivatit parcial të parë (të dy sipas x-it). Derivatet parciale të rendeve të larta njihen edhe si derivate parciale të përziera sepse ne gjejmë derivatin sipas dy ose më shumë variablave.
Konsiderojmë rastin e një funksioni me dy variabla, $f(x,y)$ pasi që të dy derivatet parciale të rendit të parë janë funksione të x dhe y, ne mund të diferencojmë secilin prej tyre sipas x ose y. Kjo do të thotë se për rastin e një funksioni me dy variabla, do ketë katër derivate të rendit të dytë.
\begin{gather*}
	\left( f_x \right)_x = f_{x x} = \pdv{}{x}\left( \pdv{f}{x} \right)= \pdv[2]{f}{x}\\
	\left( f_x \right)_y = f_{x y} = \pdv{}{y}\left( \pdv{f}{x} \right) = \pdv{f}{y}{x}\\
	\left( f_y \right)_x = f_{y x} = \pdv{}{x}\left( \pdv{f}{y} \right) = \pdv{f}{x}{y}\\
	\left( f_y \right)_y = \pdv{}{y}\left( \pdv{f}{y} \right) = \pdv[2]{f}{y}
\end{gather*}
Rasti i dytë dhe i tretë i derivateve parciale të rendit të dytë njihen edhe si derivate parciale të përziera, pasi që i diferencojmë sipas dy variabla.
\newpage

\subsection{Teorema e Klairaut-it (Clairaut's Theorem)}
Supozojmë se $f$ është i definuar në një disk D i cili përmban pikat (a,b). Nëse funksionet $f_{x y}$ dhe $f_{y x}$ janë funksione të vazhdueshme në këtë disk, atëherë,
\begin{align*}
	f_{x y}(a,b) = f_{y x}(a,b)
\end{align*}
Deri më tani i kemi shikuar vetëm derivatet parciale të rendit të dytë. Ekzistojnë edhe derivate parciale të rendeve më të larta. Më poshtë janë paraqitur disa derivate parciale të rendit të tretë të një funksioni me dy variabla.
\begin{gather*}
	f_{x y x} = \left( f_{x y} \right)_x = \pdv{}{x}\left( \pdv[2]{f}{y}{x} \right) = \frac{\partial^3 f}{\partial{x}\partial{y}\partial{x}}\\
	f_{y x x} = \left( f_{y x} \right) = \pdv{}{x}\left( \pdv[2]{f}{x}{y} \right) = \frac{\partial^3{f}}{\partial{x^2}\partial{y}}
\end{gather*}

Vërejmë se për të dyja këto, ne diferencojmë njëherë sipas y-it dhe dy herë sipas x-it. Ekziston edhe një derivat parcial i rendit të tretë në të cilin ne mund të bëjmë këtë, $f_{x x y}$. Një vazhdim i teoremës së Klairaut-it thotë se nëse trijat prej këtyre janë të vazhdueshme, atëherë ato duhet të jenë të barabarta,
\begin{align*}
	f_{x x y} = f_{x y x} = f_{y x x}
\end{align*}
Derivatet parciale e rendeve më të larta gjinden ngjashëm me këto.

\subsection{Shembull (2)}
Nga shembulli i mëparshëm [\ref{label1}] gjetëm derivatin parcial të rendit të parë, tani do të gjemë derivatin parcial të rendit të dytë.
Derivati parcial i funksionit f ishte:
\begin{align*}
	\pdv{f}{x}=4x^3 \\
	\pdv{f}{y}=\frac{3}{\sqrt{y} }
\end{align*}
Kështu që kur e diferencojmë këtë funksion prapë, do të na jep derivatin parcial të dytë
\begin{gather*}
	\pdv[2]{f}{x}= 12x^2 \\
	\pdv[2]{f}{y}= -\frac{3}{2\sqrt{y^3}}\\
	\pdv{f}{x}{y}= \pdv{f}{x}\left( \pdv{f}{y} \right) = 0
\end{gather*}
\newpage

\section{Aplikimet e derivateve parciale në matematikë}

\subsection{Minimumet dhe Maksimumet Relative}

\subsubsection{Përkufizimi}
\begin{enumerate}
	\item Një funksion $f(x,y)$ ka një minimum relativ në një pikë (a,b) nëse $f(x,y) \geq f(a,b)$ për të gjita pikat (x,y) në një zonë rreth (a,b).
	\item Një funksion f(x,y) ka një maksimum relativ në pikën (a,b) nëse $f(x,y) \leq f(a,b)$ për të gjitha pikat (x,y) në një zonë rreth (a,b).
\end{enumerate}
Vërejmë se në definicion nuk thuhet se minimumi relativ është vlera më e vogël që ai funksion do ketë. Ai thotë se vetëm në një zonë rreth pikës (a,b) funksioni do jetë më i madhë se $f(a,b)$. Jashtë asaj zone është e mundur që funksioni të jetë më i vogël. Njësoj, edhe maksimumi relativ thotë që rreth (a,b) funksioni do jetë gjithmonë më i vogël sesa f(a,b). Prapë, jashtë zonës është e mundur që funksioni të jetë më i madhë.

\subsubsection{Pikat Kritike}
Si rikujtim, pika kritike e funksionit $f(x)$ është një numër $x=c$ ashtu që $f'(c)=0$ apo $f'(c)$ nuk ekziston. Përkufizimi i pikave kritike tek funksionet me dy variabla është i ngjashëm.\\\\
Pika (a,b) është pikë kritike e $f(x,y)$ nëse njëri nga këto kushte është i vërtetë,
\begin{enumerate}
	\item$f_x(a,b)=0$ dhe $f_y(a,b) = 0$
	\item $f_x(a,b)$ dhe/ose $f_y(a,b)$ nuk ekzistojnë.
\end{enumerate}
Nëse pika (a,b) është ekstrem relativ i funksionit $f(x,y)$ dhe derivati i parë i funksionit f(x,y) ekziston tek (a,b) atëherë (a,b) është poashtu pikë kritike e $f(x,y)$.
\\\\Për ta parë këtë, le të konsiderojmë funksionin:
\begin{equation*}
	f(x,y) = xy
\end{equation*}
Dy derivatet parciale të para janë,
\begin{align*}
	f_x(x,y)=y &  & f_y(x,y)=x
\end{align*}
E vetmja pikë që do ti bëjë të dy këto derivate zero në kohën e njëjtë është (0,0) kështu që (0,0) është pikë kritike e funksionit.
\begin{center}
	\begin{figure}[H]
		\centering
		\includegraphics[width=0.7\textwidth]{grafik1.jpg}
		\caption{Grafiku i funksionit $f(x,y) = xy$ (nga WolframAlpha)}
	\end{figure}
\end{center}

\subsubsection{Testimi i pikave kritike me derivatin parcial të rendit të dytë}
Ky test përdoret për të përcaktuar nëse një pikë kritikale e funksionit është minimum lokal, apo maksimum. Supozojmë që $f(x,y)$ është një funksion real i diferencueshëm me dy variabla, derivati parcial i rendit të dytë i tij ekziston. Matrica Hessian (Hessian Matrix) H i funksionit f është matrica $2\cross 2$ e derivateve parciale të f-it,
\begin{gather}
	H(x,y) = \begin{pmatrix} f_{x x}(x,y) & f_{x y}(x,y) \\
		f_{y x}(x,y) & f_{y y}(x,y)
	\end{pmatrix}\\
	D(x,y) = det\left( H(x,y) \right) = f_{x x}(x,y)f_{y y}(x,y) - (f_{x y}(x,y))^2
\end{gather}\\
Supozojmë se (a,b) është pikë kritike e f, atëherë testi i derivatit parcial të rendit të dytë na tregon se,
\begin{enumerate}
	\item Nëse $D(a,b) > 0$ dhe $f_{x x}(a,b) > 0$ atëherë (a,b) është minimum lokal i f-it.
	\item Nëse $D(a,b) > 0$ dhe $f_{x x}(a,b) < 0$ atëherë (a,b) është maksimum lokal i f-it.
	\item Nëse $D(a,b) = 0$, kërkohen kushte shtesë.
\end{enumerate}

\subsection{Minimumet dhe Maksimumet Absolute}

\subsubsection{Përkufizimi}
\begin{enumerate}
	\item Një regjion në $\mathbb{R}^2$ thirret i mbyllur nëse ai përmban  kufinjtë e tij. Një regjion është i hapur nëse ai nuk i përfshin asnjë prej kufinjëve të tij.
	\item Një regjion në $\mathbb{R}^2$ thirret i kufizuar nëse mund të përmbahet plotësisht në një disk. Me fjalë të tjera, një rajon do të kufizohet nëse është i caktuar.
\end{enumerate}
Më poshtë janë dhënë dy definicione të drejtkëndëshit, një i hapur dhe një i mbyllur.
\begin{align*}
	\text{Hapur} &  & \text{Mbyllur}   \\
	-5 < x < 3   &  & -5 \leq x \leq 3 \\
	1 < y < 6    &  & 1 \leq y \leq 6
\end{align*}
Në rastin e parë ne nuk lejojmë që vargu të përfshijë pikat e fundit (nuk kemi përfshirë skajet e drejtkëndëshit), në fjalë të tjera ne nuk po e lejojmë që zona të përfshijë kufijtë e saj kështu që është i hapur. Në rastin e dytë ne po lejojmë që zona të përfshijë edhe kufijtë kështu që do jetë i mbyllur.
\\\\
Për ti gjetur ekstremet absolute ne duhet të:
\begin{enumerate}
	\item Gjejmë të gjitha pikat kritike te funksionit që shtrihen në zonën D dhë të gjejmë vlerën e funksionit tek ato pika.
	\item Gjejmë të gjithë ekstremet e funksionit në kufinjë.
	\item Vlerat më të mëdha dhe më të vogla që i gjejmë janë minimumi dhe maksimumi absolut i funksionit.
\end{enumerate}

\subsection{Vlerat Ekstreme të Kushtëzuara}
Vlerë ekstreme të kushtëzuar të funksionit $f(x,y)$, quajmë maksimumet dhe minimumet e tij, i cili arrihet nën kushtin që argumentet e funksioni nuk janë të pavarura ndërmjet veti, por që janë të lidhura me ekuacionin g(x,y,z)=k.

\subsubsection{Metoda e shumzuesit të Lagranzhit (Lagrange Multipliers)}
\begin{enumerate}
	\item Zgjedhim sistemin e ekuacioneve.
	      \begin{gather*}
		      \nabla f(x,y,z)=\lambda \nabla g(x,y,z)\\
		      g(x,y,z) = k
	      \end{gather*}
	\item Zëvendësojmë të gjitha zgjidhjet, (x,y,z) nga hapi i parë në $f(x,y,z)$ dhe identifikojmë vlerat minimale dhe maksimale, nëse ekzistojnë dhe $\nabla g \neq \vec{0}$ në atë pikë.
\end{enumerate}
Konstanta $\lambda$, quhet shumzues lagranzh (Lagrange Multiplier)\\
Operatori $\nabla$ quhet Nabla dhe definohet si:
\begin{gather*}
	\nabla f\equiv \pdv{f}{x} + \pdv{f}{y} + \pdv{f}{z}\\
	\nabla f = \left\langle f_x,f_y,f_z \right\rangle
\end{gather*}
\section{Aplikimet e derivateve parciale në jetën e përditshme}
Aplikimet e derivateve parciale kanë përdorim të madhë në fushën e fizikës dhe ekonomisë. Mekanika e fluideve, teoria elektromagnetike janë të modeluara me ekuacione diferenciale parciale dhe kanë plotë aplikime në jetën e përditshme.
\subsection{Ekuacioni i nxehtësisë}
Në fizikë dhe matematikë, ekuacioni i nxehtësisë është ekuacion diferencial parcial i cili përshkruan se si përhapet një madhësi (si nxehtësia) evolon sipas kohës në një medium të ngurtë, pasi spontanisht rrjedh nga vendi ku është më i lartë drejtë vendeve më të ulta. Për një funksion $u(x,y,z,t)$ me tri variabla hapsinore (x,y,z) dhe variablës së kohës t, ekuacioni i nxehtësisë është:
\begin{gather*}
	\pdv{u}{t} = \alpha\left( \pdv[2]{u}{x}+\pdv[2]{u}{y}+\dv[2]{u}{z} \right)\text{, ose,}\\
	\dot{u}=\alpha\nabla^2 u
\end{gather*}
Ku $\alpha$ është një koeficient real i quajtur difuzivitet i mediumit.
\textbf{$\nabla^2$} është operatori i Laplasit, dhe \textbf{$\dot{\textbf{u}}$} është derivati sipas kohës i \textbf{u}.
Operatori \textbf{$\nabla^2$} na jep diferencën në mes vlerës mesatare të një funksioni në rrethinën e një pike, dhe vlerën e tij në atë pikë. Kështu pra nëse \textbf{u} është temperatura, \textbf{$\nabla^2u$} na tregon nëse (dhe sa) materiali që rrethon qdo pikë është më i nxehtë apo më i ftohtë,  mesatarisht, sesa materiali në atë pikë.

\subsection{Shembull (3)}
Gjeni dimensionet e një kutie 6-fytyrëshe që ka formën e prizmit drejtkëndor me vëllimin më të madhë që mund të krijosh me $12m^2$ karton.
\\\\
Nëse përdormi të gjithë kartonin që kemi për të bërë kutinë, sipërfaqja totale A e 6-fytyrave të prizmit do jepet nga ekuacioni,
\begin{gather}
	\label{area}
	A=2xy+2yz+2zx=12
\end{gather}
Kjo është sepse prizmi drejtkëndor ka 2 fytyra xy, 2 fytyra yz dhe 2 fytyra zx. Vëllimi V i kutisë do të jetë,
\begin{gather*}
	V = xyz\\
	\text{Nëse e zgjedhim ekuacionin [\ref{area}] për z, do marrim,}\\
	z = \frac{6-xy}{x+y}\\
	\text{Zavendësojmë z në shprehjen e volumit V, dhe marrim,}\\
	V(x,y) = xy\frac{6-xy}{x+y}\\
	\text{Tani do të gjejmë pikat kritike duke gjetur derivatin parcial të parë}\\
	V_x(x,y)=-y^2\frac{x^2+2xy-6}{(x+y)^2}\\
	V_y(x,y)=-x^2\frac{y^2+2xy-6}{(x+y)^2}\end{gather*}
Tani ne duhet të zgjedhim sistemin e ekuacioneve të dhënë nga $V_x=0$ dhe $V_y=0$. Një zgjidhje triviale jepet nga pika (0,0) por nuk është e mundur fizikisht. Zgjidhjet tjera gjinden me,
\begin{gather*}
	x^2+2xy-6=0, \text{dhe,}\\
	y^2+2xy-6=0
\end{gather*}
Pas zgjidhjes do kemi,
$x=y$ dhe $x=-y$.
Zgjidha $x=-y$ nuk na duhet për këtë problem sepse të dy x dhe y janë dimensione dhe nuk mund të jenë negative. Zëvendësojmë x=y dhe kemi,
\begin{gather*}
	x^2+2x^2-6=0\\
	x=\sqrt{2}\\
	y=x=\sqrt{2} \\
	\text{Tani do të gjejmë derivatet parciale të rendit të dytë}\\
	V_{x x}(x,y) = -2y^2\frac{y^2+6}{(x+y)^3}\\
	V_{y y}(x,y) = -2x^2\frac{x^2+6}{(x+y)^3}\\
	V_{x y}(x,y) = -2xy\frac{x^2+3xy+y^2-6}{(x+y)^3}\\
	D=V_{x x}(\sqrt{2},\sqrt{2})V_{y y}(\sqrt{2},\sqrt{2})-V_{x y}^2(\sqrt{2},\sqrt{2}) = \frac{5}{2}
\end{gather*}
Pasi që, $D > 0$ dhe $V_{x x}(\sqrt{2},\sqrt{2}) < 0$ atëherë $\left( \sqrt{2},\sqrt{2}\right)$ është maksimum lokal.
\begin{gather*}
	x=\sqrt{2} \text{ [m]}\\
	y=\sqrt{2} \text{ [m]}\\
	z=\frac{6-xy}{x+y} = \sqrt{2} \text{ [m]}
\end{gather*}
\newpage

\begin{thebibliography}{10}
	\bibitem{stewart}
	James Stewart,
	\textit{Calculus: Early Transcendentals}
	Cengage Learning, 1986.

	\bibitem{paulnotes}
	Paul Dawkins, \\\texttt{http://tutorial.math.lamar.edu/Classes/CalcIII/PartialDerivsIntro.aspx}

	\bibitem{MIT}
	\text{Massachusetts Institute of Technology / Relevant Partial Derivatives articles}
	\texttt{https://ocw.mit.edu/index.htm}

	\bibitem{handbook}
	Andrei D. Polyanin,
	\textit{Handbook of Mathematics for Engineers and Scientists (Advances in Applied Mathematics)}
	Chapman and Hall/CRC, 2006
\end{thebibliography}
\end{document}
